\UseRawInputEncoding
\documentclass[12pt,a4paper]{article}

% --- Paquetes de Formato y Estilo Profesional ---
\usepackage[utf8]{inputenc}
\usepackage[spanish,es-tabla]{babel}
\usepackage[T1]{fontenc}
\usepackage{geometry}
\geometry{top=2.5cm, bottom=2.5cm, left=2.5cm, right=2.5cm}
\usepackage{bera} 
\usepackage{xcolor}
\usepackage{titlesec}
\usepackage{fancyhdr}
\usepackage{hyperref}
\usepackage{enumitem}
\usepackage{colortbl}
\usepackage{graphicx}
\usepackage{listings}
\usepackage{amsmath}
\usepackage{amssymb}
\usepackage{booktabs}
\usepackage{longtable}
\usepackage{parskip} % Para mejor manejo de parrafos y evitar paginas vacias

% --- Paleta de Colores de Ingenier\'ia ---
\definecolor{primaryblue}{RGB}{31, 73, 125}
\definecolor{secondarygray}{RGB}{120, 120, 120}
\definecolor{bgcode}{RGB}{248, 248, 248}
\definecolor{riskred}{RGB}{180, 0, 0}
\definecolor{successgreen}{RGB}{0, 120, 0}

% --- Configuración de Títulos ---
\titleformat{\section}{\color{primaryblue}\normalfont\Huge\bfseries}{\thesection}{1em}{}[\titlerule]
\titleformat{\subsection}{\color{primaryblue}\normalfont\huge\bfseries}{\thesubsection}{1em}{}
\titleformat{\subsubsection}{\color{primaryblue}\normalfont\Large\bfseries}{\thesubsubsection}{1em}{}

% --- Configuración de Bloques de Código ---
\lstset{
    basicstyle=\ttfamily\footnotesize,
    breaklines=true,
    frame=single,
    backgroundcolor=\color{bgcode},
    keywordstyle=\color{blue},
    commentstyle=\color{green!50!black},
    stringstyle=\color{red},
    tabsize=2,
    showstringspaces=false,
    captionpos=b
}

% --- Encabezado y Pie de P\'agina ---
\pagestyle{fancy}
\fancyhf{}
\lhead{\small \textbf{Proyecto Actual 24/7 - Estrategia Global}}
\rhead{\small \textbf{Sebastian Aguirre} - CTO}
\lfoot{\small \textit{Informe Técnico y Operativo Consolidado - v6.0}}
\rfoot{\textbf{P\'agina \thepage}}
\renewcommand{\headrulewidth}{0.5pt}

\begin{document}

% --- PORTADA ---
\begin{titlepage}
    \centering
    \vspace*{1.5cm}
    \Huge \textbf{PROYECTO ACUTALK 24/7} \\
    \vspace{0.8cm}
    \LARGE \textbf{Automatización Inteligente y Agregación de Repuestos para el Mercado Automotriz} \\
    \vspace{1.5cm}
    \large \textit{Especificación Detallada de Backend, Integración de Proveedores Externos y Ventajas Estratégicas de la Solución SaaS sobre Desarrollo Custom} \\
    \vspace{2.5cm}
    \textbf{Autor: Sebastian Aguirre} \\
    \small Chief Technology Officer / Arquitecto de Soluciones \\
    \vspace{3cm}
    \begin{minipage}{0.8\textwidth}
        \centering
        \textbf{Resumen:} Este documento t\'ecnico detalla la arquitectura para una red global de repuestos, centrada en la eficiencia operativa de \textbf{GoHighLevel} vinculada a \textbf{Supabase}, minimizando riesgos de mantenimiento y maximizando la conversión de ventas 24/7.
    \end{minipage}
    \vfill
    \textbf{Versión 6.0 - Documento de Ingenier\'ia Consolidado} \\
    \today
\end{titlepage}

\newpage
\tableofcontents
\newpage

\section{Introducción y Contexto del Mercado}
El mercado de autopartes se caracteriza por una fragmentación masiva de la información y una necesidad cr\'itica de precisi\'on t\'ecnica. Un error en la elección de un repuesto no solo implica una pérdida econ\'omica, sino una falla en la confianza del cliente. El proyecto \textbf{Actual 24/7} nace para solventar el cuello de botella que representa la atención humana en consultas t\'ecnicas de compatibilidad (Año, Marca, Modelo, VIN).

La visión es construir un \textbf{Hub de Suministros Inteligente} que no solo venda lo que tiene en stock, sino que actúe como un bróker automatizado consultando bases de datos propias y de terceros, entregando respuestas inmediatas con imágenes y precios competitivos en cualquier momento del día.

Esta estrategia se define como \textbf{``Autopartes Dropshipping 2.0''}, donde el valor no est\'a en el almacén físico, sino en la velocidad de la información y la capacidad de la IA para conectar una necesidad t\'ecnica (el fallo de un vehículo) con la solución log\'istica global en segundos.

\section{Infraestructura de Conectividad WhatsApp: El Punto de Acceso}
La base de la comunicación recae en la estabilidad del canal de mensajería. Para un proyecto de esta escala, no se puede utilizar una conexión común; se requiere una infraestructura de nivel empresarial.

\subsection{Especificaciones Técnicas del Número de WhatsApp}
Para garantizar que el sistema no sea bloqueado y pueda manejar altos volúmenes de mensajes, se deben cumplir los siguientes requerimientos:
\begin{itemize}
    \item \textbf{Línea Dedicada y Nueva:} Se recomienda un número de teléfono que no haya tenido reportes de spam. Debe ser una línea física o virtual con capacidad de recibir SMS para la verificación inicial de Meta.
    \item \textbf{Proceso de Calentamiento (Warm-up):} El número no puede empezar enviando 1,000 mensajes el primer día. Requiere un envío progresivo (sem 1: 50 msg/día, sem 2: 150 msg/día) para ganar reputación ante los algoritmos de WhatsApp.
    \item \textbf{Perfil de Empresa Verificado (Green Badge):} El informe contempla la gesti\'on de la verificación oficial de Meta Business Suite para aumentar la confianza del cliente y reducir la tasa de bloqueo.
    \item \textbf{Gateway de Conexión:} Uso de la API Cloud de WhatsApp (vía GHL o Whapi) que permite manejar hasta 80 mensajes por segundo, muy superior al límite de una aplicación móvil est\'andar.
\end{itemize}

\section{El Motor de Inteligencia Artificial (AI Engine)}
La "Inteligencia" del bot no reside solo en responder, sino en procesar datos t\'ecnicos complejos en milisegundos.

\subsection{Arquitectura del Modelo y Tokens}
El sistema requiere una orquestación de modelos de lenguaje (LLMs) configurados de la siguiente manera:
\begin{enumerate}
    \item \textbf{Modelo Principal (GPT-4o):} Utilizado para la identificación de repuestos y lógica de compatibilidad. Es el modelo con mayor precisi\'on t\'ecnica.
    \item \textbf{Cálculo de Tokens:} Cada consulta del cliente consume entre 500 y 1,200 tokens (incluyendo el historial de conversación y los datos t\'ecnicos de Supabase que se le envían como contexto). Esto justifica el costo mensual de la API.
    \item \textbf{Embeddings Técnicos:} Creación de una base de datos vectorial donde la IA pueda "buscar" repuestos similares por significado semántico (ej: si el cliente busca "amortiguador choto", la IA sabe que se refiere a uno de baja calidad o marca específica).
\end{enumerate}

\section{Operatividad del Backend y Servidores}
Aunque usemos GHL, el backend personalizado actúa como el ``Traductor Técnico''.

\subsection{Requerimientos del Servidor de Webhooks}
\begin{itemize}
    \item \textbf{Certificado SSL TLS 1.3:} Obligatorio para que WhatsApp y GHL confíen en enviar datos a nuestra URL.
    \item \textbf{Memoria RAM y Procesamiento:} El servidor debe ser capaz de procesar imágenes enviadas por clientes (fotos de repuestos rotos) mediante modelos de visión (Vision-LLM), lo que requiere una instancia de servidor con al menos 4GB de RAM y CPU de alta frecuencia.
    \item \textbf{Logs de Auditoría:} Registro de todas las consultas de stock externas para detectar si un proveedor est\'a fallando o enviando precios incorrectos.
\end{itemize}

\section{Requerimientos Críticos para el Inicio del Desarrollo a Medida}
Para arrancar el proyecto de forma independiente (Frontend y Backend propio), se requiere cumplir con una serie de requisitos t\'ecnicos de infraestructura que son innegociables para garantizar la estabilidad del sistema.

\subsection{1. Infraestructura de Comunicación (Número Virtual)}
\begin{itemize}
    \item \textbf{Gateway Profesional (VoIP/VAPI):} Se requiere la contratación de un proveedor de API Gateway (tipo Whapi o Meta Official Cloud API). El número debe ser una línea nueva, preferiblemente de EE.UU. o local según el mercado, para permitir el volumen de mensajes sin bloqueos.
    \item \textbf{Calentamiento del Número:} Plan de 14 días para ``humanizar'' el tráfico del número virtual antes de la automatización total.
\end{itemize}

\subsection{2. Infraestructura de Procesamiento IA}
\begin{itemize}
    \item \textbf{Acceso a Modelos de Alta Gama:} Suscripción nivel empresarial a OpenAI o Anthropic con capacidad de procesamiento de al menos 50,000 tokens por minuto (Tier 2 o superior).
    \item \textbf{Base de Datos Vectorial (Embeddings):} Se requiere un servidor de base de datos vectorial (Pinecone o Supabase Vector) para que la IA "recuerde" y busque entre miles de repuestos de forma instantánea.
\end{itemize}

\subsection{3. Conectividad y Seguridad (Webhooks)}
\begin{itemize}
    \item \textbf{Hosting del Backend:} Servidor en la nube (AWS EC2, Google Cloud o DigitalOcean) con redundancia geográfica.
    \item \textbf{Tunelización y SSL:} Conexión de Webhooks mediante certificados SSL grado bancario (TLS 1.3). Se requiere una IP est\'atica para que los Gateways de WhatsApp no rechacen la conexión.
    \item \textbf{Escalabilidad:} Configuración de contenedores (Docker) para que el servidor pueda ``crecer'' si 500 personas escriben al mismo tiempo.
\end{itemize}

\section{El Dilema del Desarrollo: Backend Custom vs. Solución SaaS}
Uno de los puntos clave del proyecto es decidir si construimos la casa desde los cimientos (Backend Custom) o si compramos un sistema modular de alta gama (\textbf{SaaS - GoHighLevel}).

\subsection{Por qué GoHighLevel (GHL) es la Solución Ganadora}
Para un proyecto de la magnitud de Actual, GHL no es solo una opción, es la estrategia correcta por las siguientes razones:
\begin{enumerate}
    \item \textbf{Cero Mantenimiento de Infraestructura:} En un backend propio (Node.js/Python), si el servidor se cae un domingo a las 3 AM, la tienda deja de vender. GHL garantiza un uptime del 99,9\% gestionado por su equipo de ingeniería global.
    \item \textbf{Gestión de Leads Integrada (CRM):} No solo estamos haciendo un chatbot. GHL nos da automáticamente el pipeline de ventas, asignación de etiquetas según la marca del carro del cliente y flujos de nutrición para que si el cliente no compró hoy, le llegue una oferta de mantenimiento en 3 meses.
    \item \textbf{Multicanalidad de serie:} Si el cliente prefiere que le enviemos la cotización por SMS o correo, GHL ya tiene estos módulos listos. Hacer esto en un backend custom requeriría integrar 3 APIs más (Twilio, SendGrid, etc.).
    \item \textbf{Velocidad de Implementación (Time-to-Market):} Desarrollar un sistema de gesti\'on de mensajes estable con manejo de archivos multimedia y cola de espera toma meses. En GHL, esto se configura en semanas mediante Workflows visuales.
    \item \textbf{Agentes Cognitivos Out-of-the-Box:} GHL permite desplegar agentes de IA que ya tienen integradas capacidades de lectura de documentos y agendamiento, eliminando la necesidad de programar toda la lógica de "entrenamiento" desde cero.
\end{enumerate}

\newpage

\section{Los Peligros del Backend Personalizado (Custom) }
Construir todo desde cero suena tentador por el "control total", pero en la industria de repuestos esto se traduce en una carga operativa inmensa.

\subsection{La Carga Amplia de Mantenimiento}
Un backend propio para Actual requeriría:
\begin{itemize}
    \item \textbf{Actualización Constante de APIs:} Proveedores externos cambian sus sistemas cada 6 meses. Un cambio mínimo en una API de un proveedor de USA rompería tu conexión, requiriendo un programador senior vigilando el código 24/7.
    \item \textbf{Gestión de Servidores y Parches:} La seguridad de los datos de pago y de clientes recaería 100\% sobre nosotros, incluyendo parches de Linux, seguridad de Node.js y mitigación de ataques DoS.
    \item \textbf{Evolución de la IA:} Los modelos de OpenAI avanzan r\'apido. Mantener la lógica de prompts sincronizada con las nuevas versiones del backend es un trabajo de tiempo completo.
\end{itemize}

\subsection{Fallos Críticos Probables}
\begin{itemize}
    \item \textbf{Latencia de APIs Externas:} Si el servidor propio debe ir a consultar a 4 proveedores en USA y China antes de responder en WhatsApp, la conexión puede colapsar o dar error de ``Timeout'', dejando al cliente con un mensaje vacío o sin respuesta.
    \item \textbf{Fugas de Memoria:} El procesamiento de imágenes pesadas (fotos de repuestos enviadas por clientes) puede saturar la RAM del servidor si no se gestiona mediante microservicios complejos, provocando caídas sistemáticas.
    \item \textbf{Pérdida de Contexto:} El fallo más común. Si el sistema no tiene una base de datos de "Estados" muy robusta, la IA podría olvidar que el cliente tiene un Mazda 3 a mitad de la conversación después de que el usuario envíe una foto.
\end{itemize}

\newpage

\section{Arquitectura del Hub de Proveedores (Supply Chain AI)}
El sistema debe centralizar múltiples fuentes de datos para ser competitivo.

\subsection{Integración vía Supabase y APIs Externas}
El backend debe actuar como un orquestador que "pregunta" de forma jerárquica:
\begin{enumerate}
    \item \textbf{Paso 1: Stock Local (Supabase):} Consulta inmediata a la tienda propia.
    \item \textbf{Paso 2: APIs Regionales:} Si no hay stock local, consulta a distribuidores nacionales con entrega el mismo día.
    \item \textbf{Paso 3: Mercado Internacional:} Consulta a catálogos de USA o China para piezas de baja rotación o alta tecnología.
\end{enumerate}

\subsection{Normalización de Datos y Sugerencias Multimedia}
La inteligencia del sistema radica en tomar datos t\'ecnicos de diferentes fuentes (que a menudo vienen en inglés o con nombres t\'ecnicos distintos) y transformarlos en una sugerencia visual clara para WhatsApp:
\begin{itemize}
    \item \textbf{Mapeador de Compatibilidad:} Si el proveedor dice "Bomba de agua para motor J20A", la IA debe saber que es compatible con la Suzuki Grand Vitara 2.0 del cliente.
    \item \textbf{Sugerencia de Imágenes:} El bot recupera automáticamente el link de la imagen. Si hay varias opciones (Original vs. Genérico), envía ambas fotos al chat para que el cliente elija visualmente.
\end{itemize}

\subsection{L\'ogica Multi-Moneda y Cálculo de Logística}
Dado que el sistema consulta proveedores en diversos países, el backend debe gestionar:
\begin{itemize}
    \item \textbf{Conversión en Tiempo Real:} Transformación instantánea de USD (USA) y CNY (China) a COP, aplicando una tasa de cambio (TRM) protegida con un margen de seguridad del 2\%.
    \item \textbf{Aranceles y Fletes:} Cálculo automático de impuestos de importación y costos de envío internacional basado en el peso volumétrico estimado de la pieza, entregando al cliente un precio final ``puesto en su puerta''.
\end{itemize}

\newpage

\section{Dise\~no Técnico de Datos (Supabase Architecture)}
La estructura de tablas es fundamental para que la búsqueda t\'ecnica sea infalible.

\begin{lstlisting}[language=SQL, caption=Esquema de Datos para Alta Disponibilidad y Compatibilidad]
-- Maestro de Productos Detallado
CREATE TABLE parts_master (
    id UUID PRIMARY KEY DEFAULT uuid_generate_v4(),
    sku VARCHAR(100) UNIQUE,
    oem_code VARCHAR(100), -- Referencia original de fabrica
    name_es TEXT,
    brand_id INT REFERENCES brands(id),
    price_cost DECIMAL(12,2),
    price_sale DECIMAL(12,2),
    stock_current INT,
    image_gallery TEXT[] -- Array de URLs para enviar por WPP
);

-- Matriz de Compatibilidad Multivehiculo
CREATE TABLE vehicle_matrix (
    part_id UUID REFERENCES parts_master(id),
    make VARCHAR(50), 
    model VARCHAR(100),
    years INT4RANGE, -- ej: [2010, 2018]
    engine_details TEXT
);
\end{lstlisting}

\newpage

\section{Ingenier\'ia de Prompts: El Cerebro del Mecánico Virtual}
Diseñar cómo responde la IA es la diferencia entre una venta y una pérdida de tiempo.

\subsection{Prompt de Sistema (System Message) - Ingenier\'ia Avanzada}
A continuación se detalla la configuración del Agente de IA mediante un prompt estructurado de alta precisi\'on, diseñado para maximizar la conversión y la exactitud t\'ecnica.

\begin{lstlisting}[caption=Instruccion de Identidad y Logica Conversacional para la IA de Actual]
[IDENTITY]
Eres el "Asesor Maestro Senior de Actual 24/7", un experto en mecanica automotriz con 20 a\~nos de experiencia y especialista en logistica de repuestos globales. Tu tono es profesional, rapido, servicial y extremadamente preciso. No eres un bot generico; eres un asistente tecnico de elite.

[CONTEXT]
Tu base de conocimientos primaria es Supabase (Stock Local) y tu red de soporte son las APIs de proveedores en USA y China. Tu objetivo es identificar la pieza exacta que el cliente necesita y cerrar la venta mediante un link de pago.

[OPERATIONAL STEPS]
1. IDENTIFICACION TECNICA: Al recibir una consulta, extrae Marca, Modelo, Anio y Pieza. Si el usuario no los proporciona, solicitalos de forma amable pero firme.
2. VALIDACION DE COMPATIBILIDAD: Antes de confirmar disponibilidad, verifica el numero VIN (Chasis) en piezas criticas (Frenos, Suspension, Motor). Si hay variacion de lado (Izquierdo/Derecho) o posicion (Delantero/Trasero), confirma con el cliente antes de cotizar.
3. CONSULTA JERARQUICA: 
   - Primero consulta la tabla 'parts_master' en Supabase.
   - Si no hay stock, informa: "Estamos consultando en nuestra red internacional..." y utiliza la API de proveedores externos.
4. ESTRATEGIA DE PRECIOS: Siempre ofrece la opcion Original (OEM) como prioridad. Si el costo es elevado, presenta una "Opcion Alternativa Certificada" (Generica de alta calidad) para asegurar la venta.

[CONSTRAINTS & RULES]
- No adivines. Si no encuentras la pieza, solicita una foto del repuesto viejo para busqueda visual.
- Respuestas breves: No escribas parrafos largos. Usa vi\~netas para precios y especificaciones.
- Formato de Imagen: Si el stock tiene URL de imagen, presentala siempre bajo el tag [IMAGE].

[CALL TO ACTION]
Al confirmar la pieza, genera confianza: "He validado la compatibilidad para tu Mazda 3 2018. Tenemos stock disponible para despacho inmediato." Finaliza con el link de pago y el tiempo estimado de entrega.
\end{lstlisting}

\newpage

\section{La Complejidad del Procesamiento Masivo de Productos}
Un punto cr\'itico para el éxito del proyecto es la gesti\'on del catálogo. Al superar la barrera de los 1,000 productos (SKUs), el procesamiento de datos deja de ser trivial y se convierte en un desafío t\'ecnico de alta complejidad.

\subsection{El Problema de la Viabilidad en Grandes Volúmenes}
\begin{itemize}
    \item \textbf{Latencia de Búsqueda Semántica:} Procesar miles de repuestos en tiempo real requiere una infraestructura de indexación muy avanzada. Si no se gestiona correctamente, el bot puede tardar más de 30 segundos en encontrar una pieza, lo que resulta en la pérdida del cliente.
    \item \textbf{Precisión del Contexto (Token Limits):} Los modelos de IA tienen un límite de ``memoria'' (context window). Al intentar que la IA elija entre 1,000 variedades de filtros de aceite, existe un alto riesgo de que entregue información incorrecta o alucine si no se usa una arquitectura RAG (Retrieval Augmented Generation) extremadamente pulida.
    \item \textbf{Limpieza de Datos:} Mantener la compatibilidad de más de 1,000 piezas requiere una labor de ingeniería de datos continua para asegurar que cada año, marca y modelo esté perfectamente mapeado.
\end{itemize}

\section{GoHighLevel (GHL): Más allá de un CRM Convencional}
GHL se selecciona como el núcleo no solo por su gesti\'on de mensajes, sino por sus capacidades avanzadas de automatización con agentes inteligentes.

\subsection{Agentes de IA y Procesamiento de Documentos (PDF)}
La capacidad de GHL para gestionar agentes inteligentes es, quizás, el factor más determinante para la viabilidad de Actual. A diferencia de un desarrollo a medida donde cada "conocimiento" debe ser programado, GHL ofrece una infraestructura de \textbf{Conversation AI} de última generación.

\subsubsection{Entrenamiento Dinámico mediante ``Knowledge Base''}
La facilidad de entrenamiento es absoluta. El sistema permite alimentar a la IA mediante la simple carga de archivos:
\begin{itemize}
    \item \textbf{Lectura Instantánea de PDFs:} Puedes subir catálogos de proveedores de 500 p\'aginas en formato PDF. El agente de GHL procesa el texto, las tablas y las referencias t\'ecnicas de forma nativa. Si un cliente pregunta por una pieza que solo aparece en ese PDF, la IA extraerá el dato con precisi\'on quirúrgica sin que un programador haya tenido que escribir una sola línea de código de búsqueda.
    \item \textbf{Sincronización de URLs:} GHL puede "leer" sitios web completos de fabricantes. Si la marca Toyota actualiza su catálogo en línea, el bot de Actual se actualiza automáticamente al rastrear esa URL.
\end{itemize}

\subsubsection{Agentes de Cierre de Ventas y Manejo de Objeciones}
No son simples bots de respuesta; son agentes comerciales entrenados para el cierre:
\begin{enumerate}
    \item \textbf{Manejo de Objeciones:} Si el cliente dice ``est\'a muy caro'', el agente de GHL puede ser configurado para ofrecer una alternativa genérica del catálogo o explicar el valor de la garantía de la pieza original de forma persuasiva.
    \item \textbf{Agendamiento y Verificación:} El agente puede verificar la disponibilidad en el calendario de instalación si el repuesto requiere montaje en taller, cerrando la cita de servicio en el mismo hilo de WhatsApp.
    \item \textbf{Triggeo de Automatizaciones:} Un agente que detecta una ``Intenci\'on de Compra'' alta puede disparar una notificación \textit{push} al celular del dueño de la tienda o crear una oportunidad en el pipeline con valor comercial estimado.
\end{enumerate}

\section{Comparativa Integral de Plataformas y Viabilidad}
Este cuadro resume por qué la dirección estrat\'egica debe apuntar hacia soluciones gestionadas.

\begin{table}[h]
\centering
\scriptsize
\begin{tabular}{|l|c|c|c|}
\hline
\rowcolor{primaryblue!10} \textbf{Criterio de Evaluación} & \textbf{Desarrollo Custom} & \textbf{GoHighLevel (GHL)} & \textbf{ManyChat} \\ \hline
\textbf{Control de Datos} & Total y Directo. & Alto (Vía Webhooks). & Limitado. \\ \hline
\textbf{Envío Multimedia (Imágenes)} & Dinámico/Programado. & Nativo y Robusto. & Estático. \\ \hline
\textbf{Carga de Mantenimiento} & \textbf{Crítica / Permanente.} & \textbf{Baja / Gestionada.} & \textbf{Muy Baja.} \\ \hline
\textbf{Riesgo de Fallos} & Alto (Código Propio). & Bajo (Infraestructura SaaS). & Mínimo. \\ \hline
\textbf{Costos de Escalamiento} & Lineales con Devs. & Bajos (Suscripción). & Medios. \\ \hline
\textbf{Gestión de Leads (CRM)} & Requiere desarrollar. & Nativo y Profesional. & Básico. \\ \hline
\rowcolor{lightgray}\textbf{EL PREFERIDO} & \textbf{Para Control Total} & \textbf{Para Ventas Masivas} & \textbf{Para SMB} \\ \hline
\end{tabular}
\caption{Matriz comparativa de infraestructura estrat\'egica.}
\end{table}

\newpage

\section{Análisis de Inversión y Modelos de Costos}
Se presentan los dos tipos de costos para el negocio: la inversión inicial (CAPEX) y el mantenimiento operativo (OPEX), detallando cada recurso necesario. Los valores en COP (Pesos Colombianos) se calculan bajo una tasa moderada y realista.

\subsection{CAPEX: Inversión Inicial (Setup y Desarrollo)}
\begin{itemize}
    \item \textbf{Número y Verificación de Línea:} \$100 - \$200 (~400,000 - 800,000 COP).
    \item \textbf{Dise\~no de Arquitectura y Webhooks:} \$2,500 - \$4,500 (~10,000,000 - 18,000,000 COP).
    \item \textbf{Normalización de Datos de Proveedores Externos:} \$1,000 (~4,000,000 COP).
    \item \textbf{Entrenamiento del Agente de IA y Embeddings:} \$800 (~3,200,000 COP).
\end{itemize}

\subsection{OPEX: Costos Mensuales Operativos (Solución GHL + Supabase)}
\begin{table}[h]
\centering
\small
\begin{tabular}{|l|p{6cm}|c|c|}
\hline
\rowcolor{primaryblue!10} \textbf{Concepto} & \textbf{Detalle} & \textbf{Costo (USD)} & \textbf{Costo Est. (COP)} \\ \hline
Número de WhatsApp & Tarificación de Meta por conversación (24h). & \$20 - \$50 & 80k - 200k \\ \hline
GoHighLevel SaaS & Licencia Agency (WhatsApp, CRM, Autom.). & \$297 & 1,188,000 \\ \hline
Supabase Cloud & Base de datos escalable (Tier Pro). & \$25 - \$50 & 100k - 200k \\ \hline
API OpenAI (Tokens) & Procesamiento de IA (GPT-4o) para 5k chats. & \$150 - \$300 & 600k - 1,200,000 \\ \hline
Gateway Whapi & Suscripción anual/mensual por canal activo. & \$30 - \$60 & 120k - 240k \\ \hline
Mantenimiento & Supervisión t\'ecnica y actualización de APIs. & \$200 & 800,000 \\ \hline
\rowcolor{gray!20} \textbf{COSTO TOTAL} & \textbf{Operación Mensual Consolidada} & \textbf{\$722 - \$1,057} & \textbf{2.9M - 4.2M COP} \\ \hline
\end{tabular}
\end{table}

\subsection{Detalle de Costos de Conectividad (Whapi / Gateways)}
Para que un sistema de autopartes sea serio, la conectividad no puede ser intermitente. El uso de \textbf{Whapi.cloud} o similares implica los siguientes costos específicos:
\begin{itemize}
    \item \textbf{Suscripción por Canal:} Whapi cobra aproximadamente \$30 USD/mes por cada canal (número de teléfono) conectado. Esto permite el envío de archivos pesados (PDFs de catálogos y fotos de alta resolución) sin las restricciones de la API oficial de Meta para cuentas que aún no son masivas.
    \item \textbf{Mantenimiento de Números Virtuales:} Si se opta por números internacionales (ej. USA), el costo de renta mensual es de aproximadamente \$5 - \$10 USD por línea.
    \item \textbf{Gestión de Proxy/IP:} Para evitar bloqueos, se recomienda el uso de un Proxy dedicado, con un costo mensual de \$10 - \$15 USD.
\end{itemize}

\subsection{Curva de Escalabilidad Financiera}
\textbf{Escenario 1: Soluci\'on Gestionada (GHL + Supabase)}
A medida que el volumen de consultas aumente de 1,000 a 10,000 repuestos procesados, los costos fijos (GHL, Supabase) se mantienen estables, mientras que los costos variables (Tokens de OpenAI y Conversaciones de Meta) crecen de forma moderada. Sin embargo, el ahorro en salarios de agentes humanos (que costarían más de 15,000,000 COP/mes para cubrir 24/7) garantiza un ROI positivo desde el segundo mes de operaci\'on.

\textbf{Escenario 2: Soluci\'on a Medida (Custom Backend)}
En una soluci\'on a medida, la curva de costos no es plana. Al escalar:
\begin{itemize}
    \item \textbf{Costo de Infraestructura:} El servidor de \$40 USD inicial tendr\'ia que escalar a instancias de \$200 - \$500 USD para manejar la concurrencia de miles de usuarios, adem\'as de costos adicionales por balanceadores de carga y CDN.
    \item \textbf{Costo de Soporte T\'ecnico:} Requiere m\'inimo dos desarrolladores \textit{on-call} para resolver ca\'idas de servidor o bugs en las integraciones de APIs que cambian constantemente. Esto representa un gasto adicional de \$4,000 - \$7,000 USD mensuales solo en n\'omina t\'ecnica.
    \item \textbf{Deuda T\'ecnica:} Cada nueva funci\'on agregada a un sistema custom aumenta la complejidad y el costo de mantenimiento a futuro, a diferencia de GHL donde las mejoras de plataforma son gratuitas y autom\'aticas.
\end{itemize}

\newpage

\newpage

\section{Seguridad de la Información y Cumplimiento}
Al manejar datos de vehículos y transiciones financieras, el sistema debe cumplir con est\'andares de seguridad elevados:
\begin{enumerate}
    \item \textbf{Protección de Datos Personales:} Almacenamiento cifrado de números de teléfono y VINs de clientes bajo normativas de Habeas Data.
    \item \textbf{Cifrado de Extremo a Extremo:} Todas las comunicaciones entre el Gateway de WhatsApp, el Backend y Supabase est\'an protegidas por certificados SSL/TLS 1.3.
    \item \textbf{Seguridad en Pagos}: Integración con pasarelas certificadas (tipo Stripe o Wompi) donde los datos de tarjeta nunca tocan nuestro servidor, reduciendo el riesgo de fraude.
\end{enumerate}

\section{Escalabilidad Futura y Optimización Post-Lanzamiento}
El proyecto no termina en el lanzamiento; se prevé una evolución constante:
\begin{itemize}
    \item \textbf{A/B Testing de Prompts:} Optimización continua de las respuestas de la IA para medir qué lenguaje genera mayor tasa de conversión.
    \item \textbf{Módulo de Dashboards Avanzados:} Visualización en tiempo real de qué repuestos son los más buscados pero no est\'an en stock, para guiar las compras de la empresa.
    \item \textbf{Integración con Marketplaces:} Capacidad de publicar automáticamente el stock de Supabase en Mercado Libre o Facebook Marketplace mediante el mismo backend.
\end{itemize}

\section{Hoja de Ruta de Implementación Ultra-Rápida (2 Semanas)}
Dado que el requerimiento es de ejecución cr\'itica e inmediata, el plan se estructura en una "Guerra de Desarrollo" (Sprint) de 14 días para la salida al mercado.

\begin{enumerate}
    \item \textbf{Semana 1: Infraestructura y Datos (Días 1-7):}
    \begin{itemize}
        \item Días 1-2: Compra de número virtual, verificación de Meta y conexión de Webhook inicial.
        \item Días 3-5: Carga masiva del inventario (1,000+ SKUs) en Supabase y normalización de activos (Fotos).
        \item Días 6-7: Configuración de los agentes de IA y entrenamiento con los PDFs de proveedores.
    \end{itemize}
    \item \textbf{Semana 2: Integración y Lanzamiento (Días 8-14):}
    \begin{itemize}
        \item Días 8-10: Desarrollo de los adaptadores para APIs de proveedores externos (USA/Nacionales).
        \item Días 11-12: Vinculación con GHL (Pipelines y CRM) y pruebas de flujo de pago.
        \item Días 13-14: Pruebas de carga final y Salida Oficial a Producción 24/7.
    \end{itemize}
\end{enumerate}

\section{Personal Técnico Requerido para Arranque a Medida}
Para ejecutar este plan en 2 semanas, el equipo debe contar con:
\begin{itemize}
    \item \textbf{1 Senior Backend Engineer:} Especialista en Node.js/Python y manejo de APIs.
    \item \textbf{1 AI Engineer:} Para la calibración de prompts, embeddings y RAG.
    \item \textbf{1 Data Architect:} Para la gesti\'on y limpieza de los 1,000+ productos en Supabase.
\end{itemize}

\newpage

\section{Auditoría Técnica del Estado Anterior y Justificación de la Nueva Arquitectura}
Es fundamental dejar constancia de por qué el trabajo realizado previamente por terceros no cumple con los est\'andares mínimos de escalabilidad y precisi\'on requeridos para el Proyecto Actual 24/7. Tras el an\'alisis del estado anterior, se detectaron las siguientes falencias cr\'iticas:

\subsection{1. Inconsistencia en el Desarrollo del Código}
Se observó una falta de progreso tangible durante un periodo de 4 meses. El código presentado inicialmente fue descartado por el desarrollador anterior a último momento, alegando un cambio total de arquitectura. Este comportamiento evidencia que prácticamente no hubo un trabajo de desarrollo sólido ni una planificación t\'ecnica real durante dicho tiempo, dejando el proyecto en un estado de incertidumbre operativa, plagado de ``deuda t\'ecnica'' y falta de documentación básica.

\subsection{2. Limitaciones del Panel Administrativo}
El panel mencionado se identifica como una ``réplica de baja calidad'' de paneles de tiendas genéricas. Sus capacidades son extremadamente limitadas y no est\'an alineadas con la complejidad de una red global de repuestos:
\begin{itemize}
    \item \textbf{Arquitectura Obsoleta:} Uso de tecnologías no escalables que colapsarían ante un tráfico de más de 50 usuarios simultáneos.
    \item \textbf{Funcionalidad Básica:} Solo incluye funciones rudimentarias como registro de ventas manual y analíticas simples que no integran el flujo conversacional.
    \item \textbf{Falta de Integración:} El panel no tiene capacidad de gestionar el mapeo OEM ni la compatibilidad din\'amica multivehículo necesaria para automatizar el negocio.
\end{itemize}

\subsection{3. Inviabilidad de los Agentes de IA Previos}
Los supuestos agentes de IA mencionados en el estado anterior carecen de robustez t\'ecnica. Se identifican como implementaciones superficiales y carentes de entrenamiento semántico RAG:
\begin{itemize}
    \item \textbf{Falta de Madurez:} La lógica de los agentes parece haber sido improvisada recientemente, careciendo de la capacidad de leer documentos t\'ecnicos (PDFs) o interactuar con bases de datos externas de forma estructurada.
    \item \textbf{Desconexión Operativa:} No presentan una integraci\'on real con CRM o sistemas de pago, lo que los reduce a simples enviadores de texto sin valor comercial.
\end{itemize}

Esta situaci\'on justifica plenamente la decisi\'on estrat\'egica de Sebastian Aguirre de pivotar hacia una arquitectura profesional basada en \textbf{GoHighLevel + Supabase}, asegurando un activo digital real y funcional para la empresa.

\newpage

\section{Conclusiones y Recomendaciones Finales}
El \textbf{Proyecto Actual 24/7} es la herramienta necesaria para que Sebastian Aguirre escale el negocio de repuestos a un nivel global. 

\textbf{El dictamen final es claro:} La implementación debe hacerse utilizando \textbf{GoHighLevel} como el núcleo de atención y CRM, vinculándolo mediante un backend ligero a \textbf{Supabase} para la gesti\'on t\'ecnica de piezas. Esto nos otorga el control de los datos pero descarga la responsabilidad de la infraestructura en plataformas SaaS de nivel mundial, reduciendo drásticamente los fallos operativos y permitiendo un crecimiento ilimitado.

\end{document}
